\documentclass[12pt]{article}
\usepackage[utf8]{inputenc}
\usepackage[spanish]{babel}
\usepackage{amsmath}
\usepackage{amsfonts}
\usepackage{amssymb}
\usepackage{geometry}
\geometry{letterpaper, margin=2.5cm}

\begin{document}

\begin{center}
    \Large \textbf{Evaluación de Examen 1: Unidad I} \\
    \large Física para Ciencias de la Salud (005-1713) \\
    \vspace{0.5cm}
    \textbf{Estudiante:} Aramerc Acuña \quad \textbf{C.I.:} 33.279.962
    \vspace{0.3cm} \\
    \large \textbf{Calificación Final: 8/10 puntos}
\end{center}

\vspace{0.5cm}

\section*{Análisis de Resultados}

\subsection*{1. Ángulo plano (Conversión a radianes)}
\textbf{Procedimiento y resultado:} Aplica correctamente la regla de tres usando la equivalencia $90^\circ = \frac{\pi}{2}$ rad. El cálculo intermedio la lleva a $\frac{37,5\pi}{90}$, obteniendo finalmente el resultado correcto de $0,41\pi$ rad (aproximación de $\frac{5\pi}{12}$). \\
\textbf{Veredicto:} \textbf{Correcto} (2/2 pts).

\subsection*{2. Sistemas de coordenadas (Longitud del hueso)}
\textbf{Procedimiento y resultado:} Construye el gráfico identificando el triángulo rectángulo y sus catetos ($a=6$, $b=8$). Plantea la fórmula del Teorema de Pitágoras ($c^2 = a^2 + b^2$) y calcula la longitud exacta de la hipotenusa $L = 10$ cm. Excelente procedimiento. \\
\textbf{Veredicto:} \textbf{Correcto} (2/2 pts).

\subsection*{3. Vectores (Fuerza resultante)}
\textbf{Procedimiento y resultado:} Entiende el concepto de suma vectorial agrupando las componentes $\hat{i}$ y $\hat{j}$. Sin embargo, al igual que en casos anteriores, comete un \textbf{error de transcripción y de signo} en la componente $\hat{j}$: en lugar de sumar $(25 + 35)\hat{j}$, escribe $(25 - 15)\hat{j}$ (aparentemente arrastrando el 15 de la componente $\hat{i}$). Esto la lleva a un vector incorrecto de $30\hat{i} - 10\hat{j}$ N. \\
\textbf{Veredicto:} \textbf{Parcialmente correcto} (Planteamiento válido, error aritmético/transcripción al operar) (1/2 pts).

\subsection*{4. Potencias de diez (Notación científica y prefijo)}
\textbf{Procedimiento y resultado:} Convierte la cifra decimal adecuadamente a una expresión con potencias de diez ($12,5 \times 10^{-6}$). Al igual que sus compañeros, omite la segunda parte de la pregunta: indicar el prefijo adecuado del Sistema Internacional (micrómetros, $\mu$m). \\
\textbf{Veredicto:} \textbf{Incompleto / Parcialmente correcto} (1/2 pts).

\subsection*{5. Sistemas de unidades (Conversión de densidad)}
\textbf{Procedimiento y resultado:} Plantea factores de conversión formales de manera errónea, pero luego ejecuta el procedimiento aritmético correcto por separado: divide entre 1000 para la masa y multiplica por $1.000.000$ para el volumen. Obtiene el valor numérico exacto ($1030$), aunque olvida colocar las unidades ($\text{kg/m}^3$) en su respuesta final remarcada. Se le otorga la puntuación completa por el cálculo numérico exacto. \\
\textbf{Veredicto:} \textbf{Correcto} (2/2 pts).

\vspace{0.5cm}
\section*{Comentarios Generales y Sugerencias}
La estudiante muestra un buen dominio general de los conceptos básicos de geometría y conversión aritmética. Se sugieren los siguientes aspectos a mejorar:
\begin{itemize}
    \item \textbf{Cuidado con los signos y transcripción:} Revisar cuidadosamente los datos (como en la Pregunta 3) antes de realizar operaciones algebraicas.
    \item \textbf{Uso de unidades:} En física, un resultado numérico carece de significado completo sin sus unidades. Recordar siempre incluir las unidades en el resultado final (como faltó en la Pregunta 5).
    \item \textbf{Leer los enunciados completos:} Asegurarse de responder todas las partes de la pregunta, como se observó en la Pregunta 4.
\end{itemize}

\end{document}